% ═══════════════════════════════════════════════════════════════
%  ch07-ring-buffer.tex
%  Кольцо кадров: как не потерять данные
% ═══════════════════════════════════════════════════════════════

\chapter{Кольцо кадров}
\label{ch:ring-buffer}

\epigraph{<<Данные приходят быстрее, чем уходят. Кольцо~---
это способ не потерять их по дороге>>.}{}

\section{Введение}
\label{sec:ring-intro}

В \texttt{ch05} мы настроили ADC и DMA. В \texttt{ch06}~---
добавили триггер. Теперь \emph{данные} \emph{текут}
\emph{непрерывно}. Но \emph{куда} их \emph{девать}?

\textbf{Проблема}: DMA \emph{пишет} с \emph{частотой}
\SI{100}{\kilo\hertz}. USB CDC \emph{отправляет} с
\emph{частотой} \SI{1}{\mega\hertz} (\emph{по} \emph{шине}),
но \emph{пакетами}~--- \emph{не} \emph{мгновенно}.
Пока USB \emph{отправляет} \emph{один} кадр~--- DMA
\emph{уже} \emph{заполнил} \emph{следующий}.

\textbf{Где} \emph{хранить} \emph{готовые} кадры, пока
USB \emph{их} \emph{отправляет}?

\begin{ideabox}
\textbf{Аналогия: конвейер.} На \emph{заводе} \emph{детали}
\emph{едут} по \emph{конвейеру}. \emph{Рабочий} (DMA)
\emph{ставит} \emph{новую} деталь. \emph{Упаковщик} (USB)
\emph{берёт} \emph{готовую}. \emph{Между} ними~---
\emph{накопитель} (кольцо). Если \emph{рабочий} \emph{быстрее}
\emph{упаковщика}~--- детали \emph{копятся}. Если \emph{накопитель}
\emph{полон}~--- \emph{рабочий} \emph{ждёт}.
\end{ideabox}

\subsection{Что вы узнаете}

\begin{itemize}
    \item Что такое \emph{кольцевой буфер} и зачем он нужен.
    \item Как \emph{синхронизировать} \emph{писателя} (DMA)
          и \emph{читателя} (USB).
    \item Что делать при \emph{переполнении}.
    \item Как \emph{реализовать} кольцо кадров на STM32.
\end{itemize}

\section{Кольцевой буфер}
\label{sec:ring-what}

\textbf{Кольцевой буфер} (ring buffer, circular buffer)~---
это \emph{массив} \emph{фиксированного} размера, который
\emph{работает} \emph{по} \emph{кругу}.

\textbf{Как} это \emph{работает}:
\begin{enumerate}
    \item \emph{Писатель} \emph{кладёт} данные \emph{в}
          ячейку \texttt{write\_idx}.
    \item \emph{Увеличивает} \texttt{write\_idx} \emph{на}
          1 (\emph{по} \emph{модулю} \emph{размера}).
    \item \emph{Читатель} \emph{берёт} данные \emph{из}
          ячейки \texttt{read\_idx}.
    \item \emph{Увеличивает} \texttt{read\_idx} \emph{на} 1.
    \item \emph{Когда} \emph{дошли} \emph{до} \emph{конца}~---
          \emph{возвращаемся} \emph{в} \emph{начало}.
\end{enumerate}

\begin{figure}[htbp]
\centering
\adjustbox{max width=\textwidth}{%
\begin{tikzpicture}[
    x=1cm, y=1cm,
    cell/.style={
        draw, thick,
        minimum width=1.5cm,
        minimum height=1.2cm,
        align=center,
        font=\footnotesize\ttfamily
    },
    full/.style={cell, fill=clrLogicData!30},
    empty/.style={cell, fill=white},
    write/.style={cell, fill=clrWarn!30},
    read/.style={cell, fill=clrUIData!30}
]

% ─── Координатная сетка ────────────────────────────
\draw[step=1cm, gray!20, very thin] (-1, -2) grid (11, 2);

% Ячейки кольца
\node[read]  (c0) at (0, 0) {Кадр 0};
\node[read]  (c1) at (2, 0) {Кадр 1};
\node[full]  (c2) at (4, 0) {Кадр 2};
\node[full]  (c3) at (6, 0) {Кадр 3};
\node[write] (c4) at (8, 0) {Кадр 4};
\node[empty] (c5) at (10, 0) {пусто};

% Индексы
\draw[-{Latex[length=3mm]}, very thick, color=clrUIData]
    (c0.south) -- (0, -1) node[below] {\texttt{read\_idx}};
\draw[-{Latex[length=3mm]}, very thick, color=clrWarn]
    (c4.south) -- (8, -1) node[below] {\texttt{write\_idx}};

% Подписи
\node[font=\footnotesize, above] at (c0.north) {читается};
\node[font=\footnotesize, above] at (c4.north) {пишется};

% Метки осей (для отладки)
\foreach \x in {-1,...,11} {
    \node[font=\tiny, gray] at (\x, -2.3) {\x};
}
\foreach \y in {-2,-1,0,1,2} {
    \node[font=\tiny, gray] at (-1.3, \y) {\y};
}

\end{tikzpicture}
}% end adjustbox
\caption{Кольцевой буфер из шести кадров. \texttt{read\_idx}
указывает на текущий читаемый кадр. \texttt{write\_idx}~---
на текущий записываемый.}
\label{fig:ring-buffer}
\end{figure}

\begin{keybox}
\textbf{Почему кольцо?} Потому что \emph{память}
\emph{ограничена}. Мы \emph{не} можем \emph{хранить}
\emph{все} кадры. Но \emph{можем} \emph{хранить}
\emph{последние} \texttt{N}. \emph{Старые}~---
\emph{перезаписываются}.
\end{keybox}

\section{Синхронизация: писатель и читатель}
\label{sec:ring-sync}

\textbf{Проблема}: \emph{писатель} (DMA) и \emph{читатель}
(USB)~--- \emph{разные} \emph{сущности}. Они \emph{работают}
\emph{одновременно}. \emph{Что} если \emph{писатель}
\emph{обгонит} \emph{читателя}?

\textbf{Три} \emph{состояния}:
\begin{enumerate}
    \item \textbf{Пусто} (\texttt{read\_idx == write\_idx}).
          \emph{Читатель} \emph{ждёт}.
    \item \textbf{Заполнено} (\texttt{write\_idx + 1 == read\_idx}).
          \emph{Писатель} \emph{ждёт} (\emph{некуда} писать).
    \item \textbf{Есть данные} (\emph{иначе}). Оба
          \emph{работают}.
\end{enumerate}

\begin{warnbox}
\textbf{Что если \emph{писатель} \emph{обогнал}?}
Он \emph{перезапишет} \emph{данные}, \emph{которые}
\emph{ещё} \emph{не} \emph{прочитаны}. \emph{Потеря}.
\textbf{Решение}: \emph{писатель} \emph{ждёт}, \emph{если}
\emph{кольцо} \emph{полное} (\emph{или} \emph{увеличить}
размер).
\end{warnbox}

\subsection{Атомарность}

\textbf{Как} \emph{сделать} так, \emph{чтобы} \emph{не}
\emph{было} \emph{состояния} \emph{гонки}? \emph{Ведь}
\emph{писатель} и \emph{читатель}~--- \emph{разные}
\emph{потоки} (\emph{или} \emph{прерывание} \emph{и}
\emph{основной} \emph{цикл}).

\textbf{Решение}: \emph{атомарные} \emph{операции}.
\emph{В} STM32~--- \emph{можно} \emph{запретить}
\emph{прерывания} \emph{на} \emph{короткое} \emph{время}:
\begin{listing}[htbp]
\begin{minted}[frame=none, bgcolor=gray!5]{c}
__disable_irq();
// критическая секция: обновляем индексы
__enable_irq();
\end{minted}
\caption{Критическая секция}
\label{lst:ring-critical}
\end{listing}

\begin{keybox}
\textbf{Критическая секция~--- \emph{короткая}.}
\emph{Запрет} \emph{прерываний}~--- \emph{на} \emph{несколько}
\emph{тактов}. Если \emph{дольше}~--- \emph{пропускаются}
\emph{прерывания} (ADC, USB). \textbf{Правило}:
\emph{внутри}~--- \emph{только} \emph{индексы}.
\end{keybox}

\section{Переполнение}
\label{sec:ring-overflow}

\textbf{Что} \emph{делать}, \emph{если} \emph{USB}
\emph{отстаёт}? \emph{Три} \emph{стратегии}:

\subsection{Стратегия 1: блокировка}

\textbf{Писатель} (DMA) \emph{ждёт}, \emph{пока}
\emph{читатель} \emph{не} \emph{освободит} \emph{место}.
\textbf{Плюс}: \emph{данные} \emph{не} \emph{теряются}.
\textbf{Минус}: \emph{ADC} \emph{останавливается} (\emph{плохо}).

\subsection{Стратегия 2: перезапись}

\textbf{Писатель} \emph{перезаписывает} \emph{старые}
\emph{данные}. \textbf{Плюс}: \emph{ADC} \emph{работает}.
\textbf{Минус}: \emph{данные} \emph{теряются}.

\subsection{Стратегия 3: счётчик потерь}

\textbf{Писатель} \emph{перезаписывает}, \emph{но}
\emph{увеличивает} \emph{счётчик} \emph{потерь}.
\emph{Пользователь} \emph{видит}: <<\emph{потеряно}
\emph{N} \emph{кадров}>>. \textbf{Плюс}: \emph{честно}.
\textbf{Минус}: \emph{данные} \emph{теряются}.

\begin{ideabox}
\textbf{Что выбрать?} Для \emph{осциллографа}~---
\emph{стратегия} 3. \emph{Честно} \emph{сказать}
\emph{пользователю}: <<\emph{данные} \emph{потеряны}>>~---
\emph{лучше}, \emph{чем} \emph{молча} \emph{испортить}
\emph{картинку}.
\end{ideabox}

\section{Код: ring\_buffer.c}
\label{sec:ring-code}

\textbf{Теперь}~--- \emph{код}. Он \emph{короткий}, но
\emph{плотный}. \emph{Каждая} строка \emph{важна}.

\subsection{Структура}

\begin{listing}[htbp]
\begin{minted}[frame=none, bgcolor=gray!5]{c}
#define RING_SIZE 8

typedef struct {
    osc_frame_t frames[RING_SIZE];
    volatile uint16_t write_idx;
    volatile uint16_t read_idx;
    volatile uint32_t dropped;
} RingBuffer;
\end{minted}
\caption{Структура \texttt{RingBuffer}}
\label{lst:ring-struct}
\end{listing}

\textbf{Что здесь важно}:
\begin{itemize}
    \item \texttt{frames[RING\_SIZE]}~--- \emph{массив}
          кадров.
    \item \texttt{write\_idx}~--- \emph{куда} \emph{писать}.
    \item \texttt{read\_idx}~--- \emph{откуда} \emph{читать}.
    \item \texttt{dropped}~--- \emph{счётчик} \emph{потерь}.
    \item \texttt{volatile}~--- \emph{чтобы} \emph{компилятор}
          \emph{не} \emph{оптимизировал}.
\end{itemize}

\subsection{Запись}

\begin{listing}[htbp]
\begin{minted}[frame=none, bgcolor=gray!5]{c}
bool ring_push(RingBuffer *rb, const osc_frame_t *frame)
{
    uint16_t next = (rb->write_idx + 1) % RING_SIZE;

    if (next == rb->read_idx) {
        // Кольцо полное
        rb->dropped++;
        return false;
    }

    rb->frames[rb->write_idx] = *frame;
    __disable_irq();
    rb->write_idx = next;
    __enable_irq();
    return true;
}
\end{minted}
\caption{\texttt{ring\_push()}~--- запись}
\label{lst:ring-push}
\end{listing}

\textbf{Что здесь важно}:
\begin{itemize}
    \item \texttt{next}~--- \emph{следующий} \emph{индекс}
          (\emph{по} \emph{кругу}).
    \item \texttt{next == read\_idx}~--- \emph{кольцо}
          \emph{полное}.
    \item \texttt{dropped++}~--- \emph{увеличиваем}
          \emph{счётчик}.
    \item \texttt{*frame}~--- \emph{копируем} \emph{кадр}
          \emph{целиком}.
    \item \texttt{\_\_disable\_irq()}~--- \emph{атомарно}
          \emph{обновляем} \emph{индекс}.
\end{itemize}

\subsection{Чтение}

\begin{listing}[htbp]
\begin{minted}[frame=none, bgcolor=gray!5]{c}
bool ring_pop(RingBuffer *rb, osc_frame_t *out)
{
    if (rb->read_idx == rb->write_idx) {
        return false;  // пусто
    }

    *out = rb->frames[rb->read_idx];
    __disable_irq();
    rb->read_idx = (rb->read_idx + 1) % RING_SIZE;
    __enable_irq();
    return true;
}
\end{minted}
\caption{\texttt{ring\_pop()}~--- чтение}
\label{lst:ring-pop}
\end{listing}

\textbf{Что здесь важно}:
\begin{itemize}
    \item \texttt{read\_idx == write\_idx}~--- \emph{пусто}.
    \item \texttt{*out = ...}~--- \emph{копируем} \emph{кадр}
          \emph{целиком}.
    \item \texttt{read\_idx = ...}~--- \emph{сдвигаем}
          \emph{индекс}.
\end{itemize}

\begin{warnbox}
\textbf{Копируйте \emph{целиком}, а не \emph{по}
\emph{указателю}.} Писатель \emph{может} \emph{перезаписать}
кадр, \emph{пока} читатель \emph{его} \emph{читает}.
\emph{Копия}~--- \emph{безопаснее}.
\end{warnbox}

\section{Уроки главы}
\label{sec:ring-lessons}

\begin{enumerate}
    \item \textbf{Кольцо~--- это \emph{накопитель}.} Данные
          \emph{копятся} \emph{между} DMA и USB.
    \item \textbf{Три состояния}: \emph{пусто}, \emph{полно},
          \emph{есть данные}.
    \item \textbf{Атомарность}: \emph{индексы} \emph{обновляем}
          \emph{в} \emph{критической} \emph{секции}.
    \item \textbf{Переполнение}: \emph{три} \emph{стратегии}.
          Для \emph{осциллографа}~--- \emph{счётчик} потерь.
    \item \textbf{Копирование}: \emph{копируйте} \emph{кадр}
          \emph{целиком}.
    \item \textbf{\texttt{volatile}}: \emph{обязательно}
          для \emph{индексов}.
\end{enumerate}

\section{Чек-лист грабель}
\label{sec:ring-pitfalls}

\begin{table}[htbp]
\centering
\caption{Грабли кольцевого буфера}
\label{tab:ring-pitfalls}
\small
\begin{tabularx}{\textwidth}{@{}c L L@{}}
\toprule
\textbf{№} & \textbf{Грабля} & \textbf{Решение} \\
\midrule
1 & Нет \texttt{volatile} & Компилятор \emph{оптимизирует} \\
2 & Нет критической секции & Гонка данных \\
3 & Передача \emph{указателя} & \emph{Копия} \\
4 & Нет счётчика потерь & \emph{Честность} \\
5 & \emph{Слишком} маленькое кольцо & Увеличить \\
6 & \emph{Долгая} критическая секция & Только индексы \\
7 & Нет проверки на \emph{пусто} & \emph{Падение} \\
8 & Нет проверки на \emph{полно} & \emph{Потеря} \\
\bottomrule
\end{tabularx}
\end{table}

\section{Упражнения}
\label{sec:ring-exercises}

\begin{exercisebox}
\textbf{Упражнение 1.} Нарисуйте кольцо из \num{4} кадров.
Покажите \texttt{read\_idx} и \texttt{write\_idx} для
\emph{трёх} случаев: пусто, полно, два кадра.
\end{exercisebox}

\begin{exercisebox}
\textbf{Упражнение 2.} Что произойдёт, если \emph{убрать}
\texttt{\_\_disable\_irq()} из \texttt{ring\_push()}?
Приведите сценарий с \emph{гонкой}.
\end{exercisebox}

\needspace{5\baselineskip}
\begin{exercisebox}
\textbf{Упражнение 3.} Реализуйте \emph{стратегию}
\emph{блокировки} (\emph{ждать}, \emph{пока} \emph{место}
\emph{освободится}). Как это \emph{повлияет} на ADC?
\end{exercisebox}

\begin{exercisebox}
\textbf{Упражнение 4.} Предложите способ \emph{передачи}
\emph{указателя} вместо \emph{копирования} \emph{кадра}.
Как \emph{избежать} гонки?
\end{exercisebox}

\section{Проверьте сами}
\label{sec:ring-check}

\begin{enumerate}
    \item Подключите Nucleo к компьютеру.
    \item Соберите прошивку (см. главу 2).
    \item Подключите генератор к \texttt{PA0} (\SI{1}{\kilo\hertz},
          \SI{1}{\volt}pp).
    \item Запустите OpenOCD + GDB.
    \item Поставьте точку останова в \texttt{ring\_push()}.
    \item Посмотрите, как \emph{меняются} индексы.
    \item Сделайте \emph{скриншот} — \emph{пришлите} в книгу.
\end{enumerate}

\begin{ideabox}
\textbf{Если получилось~--- вы поняли кольцо.} Это
\emph{главное}. Без него~--- \emph{потери} данных.
\emph{Есть} \emph{только} \emph{обрывки} сигнала. Дальше~---
\emph{USB CDC}, \emph{графика}, \emph{осциллограмма}.
\emph{Всё} \emph{впереди}.
\end{ideabox}
